\begin{tabularx}{\textwidth}{L{35mm}X}
\toprule
\textbf{Elvárás vagy tétel} & \textbf{Lényege} \\
\midrule
Veszteségmentes felbontás & A felbontott relációk természetes joinjából visszaállítható az eredeti reláció, vagyis nem keletkezik hamis rekord és nem vész el információ. \\
Függőségmegőrzés & Az eredeti funkcionális függőségek ellenőrizhetők a felbontott sémákon is anélkül, hogy mindig össze kellene joinolni őket. \\
Heath-tétel & Ha \(R(A,B,C)\) relációban \(A \to B\) teljesül, akkor az \(R_1(A,B)\) és \(R_2(A,C)\) felbontás veszteségmentes. \\
Rissanen-tétel & Egy felbontás akkor tekinthető függetlennek, ha a függőségrendszer előállítható a részsémák függőségeiből, és a közös attribútum megfelelő kulcsszerepet tölt be valamelyik részsémában. \\
Normalizálási lépés & A redundanciát okozó függőség bal oldalát és a tőle függő mezőket külön relációba tesszük, majd PK--FK kapcsolattal kötjük össze a részeket. \\
Kompromisszum & A normalizálás csökkenti az anomáliákat, de több joinhoz vezethet; teljesítmény miatt néha szabályozott denormalizálás indokolt. \\
\bottomrule
\end{tabularx}
